\begin{proof}
\textbf{第一步：将条件尾平方函数识别为角块的模。}
令
\[
 R_n=X-\mathbb E_nX,\qquad
 B_n=(I-P_n)XP_n,\qquad A_n=P_nX(I-P_n).
\]
壳条件期望的定义给出
\[
 R_nP_n=B_n,\qquad P_nR_n=A_n,
 \qquad P_n\mathbb E_n(a)P_n=P_naP_n.
\]
因此
\begin{equation}\label{sch:eq:tail-moduli}
 s_{n,c}(X)^2=B_n^*B_n,\qquad
 s_{n,r}(X)^2=A_nA_n^*.
\end{equation}
特别地，
\[
 \norm{s_{n,c}(X)}_p=\norm{B_n}_p,
 \qquad \norm{s_{n,r}(X)}_p=\norm{A_n}_p.
\]
再由 \eqref{sch:eq:difference}，当 $k>n$ 时，
\[
 \Delta_kX\,P_n=E_kXP_n,\qquad
 P_n\Delta_kX=P_nXE_k.
\]
利用 $\sum_{k>n}E_k=I-P_n$，立即得到
\eqref{sch:eq:tail-diff-column}--\eqref{sch:eq:tail-diff-row}。
这些正部分和单调递增、有界且支撑在有限维的 $P_nH$ 上，故按算子范数收敛。

\medskip
\textbf{第二步：按块列作酉分解。}
对每个 $n$，取 $E_nH$ 的一组标准正交基
$\{\eta_{n,j}:1\leq j\leq d_n\}$。
所有这些向量合起来构成 $H$ 的标准正交基。映射
\[
 U:\HS\longrightarrow\bigoplus_{n\geq1}\bigoplus_{j=1}^{d_n}H,
 \qquad U(A)=(A\eta_{n,j})_{n,j}
\]
是酉算子，因为
\[
 \norm A_2^2=\sum_{n\geq1}\sum_{j=1}^{d_n}\norm{A\eta_{n,j}}_H^2,
\]
且任意平方可和的向量列均唯一确定一个 Hilbert--Schmidt 算子。
由于 $\PP$ 是 \eqref{eq:martingale-transform} 的三角鞅变换，
对其第 $n$ 个块列有
\[
 \PP(A)\eta_{n,j}=P_nA\eta_{n,j}.
\]
从而
\[
 \CC_X(A)\eta_{n,j}=[X,P_n]A\eta_{n,j},
\]
即
\begin{equation}\label{sch:eq:unitary}
 \boxed{\quad
 U\CC_XU^{-1}
 =\bigoplus_{n\geq1}\bigoplus_{j=1}^{d_n}[X,P_n].
 \quad}
\end{equation}
相对于 $H=P_nH\oplus(I-P_n)H$，有
\[
 K_n=[X,P_n]=\begin{pmatrix}0&-A_n\\B_n&0\end{pmatrix},
 \qquad
 K_n^*K_n=\begin{pmatrix}B_n^*B_n&0\\0&A_n^*A_n\end{pmatrix}.
\]
于是每个 $K_n$ 都是有限秩，且
\begin{equation}\label{sch:eq:local-schatten}
 \norm{K_n}_p^p
 =\norm{B_n}_p^p+\norm{A_n}_p^p
 =\norm{s_{n,c}(X)}_p^p+\norm{s_{n,r}(X)}_p^p.
\end{equation}

\medskip
\textbf{第三步：汇集各分量的奇异值。}
若 \eqref{sch:eq:criterion} 右侧级数收敛，则由
\eqref{sch:eq:local-schatten} 有 $\norm{K_n}^p\leq\norm{K_n}_p^p\to0$。
故 \eqref{sch:eq:unitary} 的直和是紧算子；
各个 $K_n$ 的非零奇异值在相应分量中重复 $d_n$ 次。
求 $p$ 次幂和即得 \eqref{sch:eq:main} 及 \eqref{sch:eq:corners}。
反之，若 $\CC_X\in S_p(\HS)$，则有限分量直和的奇异值 $p$ 次幂和一致有界；
对分量数取极限并使用 \eqref{sch:eq:local-schatten}，即得
\eqref{sch:eq:criterion} 右侧级数收敛及同一精确等式。
这对每个 $p>0$ 都成立。
\end{proof}

